\section{Dữ liệu và benchmark}
\label{sec:data}

\subsection{Nguồn dữ liệu pháp luật}

Corpus được xây dựng từ nội dung của Bộ luật Hình sự Việt Nam năm 2015 và
các nội dung sửa đổi, bổ sung có liên quan
\cite{quochoi2015blhs,quochoi2017suadoiblhs}. Nguồn dữ liệu tập trung vào
các quy định mô tả điều kiện áp dụng và hệ quả pháp lý trong luật hình sự.
Nghiên cứu không sử dụng bản án, hồ sơ vụ án, dữ liệu cá nhân hoặc dữ liệu
hành vi thực tế.

Đơn vị đầu vào của pipeline không phải toàn bộ điều luật dưới dạng một đoạn
văn duy nhất, mà là các quy tắc đã được phân tách và chuẩn hóa. Một điều luật
có thể tạo ra nhiều rule nếu nó chứa nhiều cặp điều kiện--hệ quả khác nhau.
Ngược lại, nhiều điều luật có thể dẫn đến cùng một effect event nhưng áp
dụng cho các điều kiện hoặc chủ thể khác nhau.

Mỗi rule record chứa tối thiểu các trường:

\begin{itemize}
    \item \texttt{rule\_id}: mã định danh duy nhất của quy tắc;
    \item \texttt{article\_id}: số điều luật;
    \item \texttt{article\_title}: tiêu đề điều luật;
    \item \texttt{legal\_subject}: chủ thể pháp luật;
    \item \texttt{condition}: điều kiện áp dụng;
    \item \texttt{effect}: hệ quả pháp lý;
    \item \texttt{condition\_event}: mã event điều kiện;
    \item \texttt{condition\_event\_name}: tên đọc được của event điều kiện;
    \item \texttt{effect\_event}: mã event hệ quả;
    \item \texttt{effect\_event\_name}: tên đọc được của event hệ quả;
    \item \texttt{causal\_type}: loại quan hệ hoặc loại hệ quả nếu có.
\end{itemize}

Corpus pháp luật được sử dụng như một nguồn tri thức quy phạm. Hệ thống
không suy ra rằng các quan hệ điều kiện--hệ quả trong corpus phản ánh quan
hệ nhân quả thống kê ngoài thực tế. Mọi kết luận của pipeline đều bị giới
hạn bởi nội dung rule đã được mã hóa.

\subsection{Quy trình chuẩn hóa quy tắc}

Quá trình chuẩn hóa nhằm chuyển văn bản pháp luật thành một schema thống
nhất cho graph construction, vector indexing và evaluation. Quy trình gồm
các bước chính sau:

\begin{enumerate}
    \item xác định điều luật và tiêu đề tương ứng;

    \item tách nội dung thành điều kiện áp dụng và hệ quả pháp lý;

    \item xác định chủ thể chịu tác động của quy tắc;

    \item chuẩn hóa condition và effect thành các event name ngắn gọn;

    \item tạo event ID ổn định từ event name;

    \item gán rule ID và liên kết rule với article ID;

    \item kiểm tra các trường bắt buộc và loại bỏ record không đủ thông tin.
\end{enumerate}

Event ID được chuẩn hóa bằng cách chuyển nội dung sang chữ thường, bỏ dấu
tiếng Việt, thay khoảng trắng bằng dấu gạch dưới và loại bỏ ký tự không cần
thiết. Ví dụ:

\begin{quote}
Event name: ``Phạm tội chưa đạt'';

Event ID: \texttt{pham\_toi\_chua\_dat}.
\end{quote}

Các cách diễn đạt pháp luật khác nhau nhưng có cùng ý nghĩa có thể được ánh
xạ về cùng một event ID. Ngược lại, những hệ quả có khác biệt đáng kể về mức
phạt, đối tượng hoặc điều kiện áp dụng được giữ thành các event riêng biệt.
Việc chuẩn hóa này giúp giảm trùng lặp bề mặt nhưng vẫn bảo toàn sự khác biệt
cần thiết cho suy luận pháp luật.

Quy trình hiện tại không thực hiện chuẩn hóa đầy đủ đối với mọi ngoại lệ,
phủ định, tham chiếu chéo hoặc cấu trúc nhiều điều kiện. Phần lớn rule trong
corpus thực nghiệm được biểu diễn dưới dạng một condition event dương và một
effect event dương. Do đó, corpus phù hợp nhất với suy luận chuỗi điều
kiện--hệ quả ngắn.

\subsection{Thống kê corpus và đồ thị}

Corpus sau chuẩn hóa gồm \(2\,884\) rule và \(1\,814\) event. Đồ thị dị thể
được xây dựng từ corpus chứa \(5\,134\) node và \(14\,420\) edge. Các thống
kê chính được trình bày trong Bảng~\ref{tab:graph}.

\begin{table}[t]
    \centering
    \caption{Thống kê corpus và đồ thị nhân quả pháp luật.}
    \label{tab:graph}
    \begin{tabular}{lr}
        \hline
        \textbf{Thành phần} & \textbf{Số lượng} \\
        \hline
        Quy tắc pháp luật & \(2\,884\) \\
        Sự kiện pháp lý & \(1\,814\) \\
        Tổng số node & \(5\,134\) \\
        Tổng số edge & \(14\,420\) \\
        Cạnh \texttt{CAUSES} & \(2\,884\) \\
        Cặp event có hướng duy nhất & \(2\,122\) \\
        Chuỗi nhân quả hai bước & \(112\) \\
        \hline
    \end{tabular}
\end{table}

Số cạnh \texttt{CAUSES} bằng số rule vì mỗi rule tạo ít nhất một quan hệ từ
condition event đến effect event. Tuy nhiên, số cặp event duy nhất nhỏ hơn
số cạnh \texttt{CAUSES}, cho thấy một số cặp event được hỗ trợ bởi nhiều rule
hoặc nhiều căn cứ pháp luật khác nhau.

Tổng số edge lớn hơn đáng kể so với số cạnh \texttt{CAUSES} vì graph còn lưu
các liên kết dị thể giữa event, rule và article. Các edge này hỗ trợ truy
vết nguồn evidence nhưng không được xem là causal edge khi thực hiện path
traversal.

Từ projection event--event, hệ thống xác định \(112\) chuỗi nhân quả hai
bước có dạng:

\[
v_0 \rightarrow v_1 \rightarrow v_2.
\]

Một chuỗi có thể được sử dụng để tạo nhiều câu hỏi với các answer intent
khác nhau, chẳng hạn hỏi outcome, hỏi nguyên nhân, hỏi mediator hoặc yêu cầu
giải thích toàn bộ path. Vì vậy, số câu hỏi benchmark có thể lớn hơn số
chuỗi hai bước duy nhất.

\subsection{Mục tiêu xây dựng benchmark}

Benchmark được xây dựng nhằm đánh giá riêng biệt các năng lực sau:

\begin{enumerate}
    \item truy hồi đúng rule và article;
    \item truy hồi đúng event;
    \item khôi phục causal path theo đúng thứ tự;
    \item phân biệt quan hệ trực tiếp và quan hệ qua mediator;
    \item xác minh claim theo nội dung câu hỏi;
    \item sinh câu trả lời đúng answer intent;
    \item trích dẫn đúng căn cứ pháp luật.
\end{enumerate}

Thiết kế benchmark nhấn mạnh việc đánh giá supporting path thay vì chỉ đánh
giá câu trả lời cuối. Cách tiếp cận này tương đồng với nguyên tắc đánh giá
supporting facts trong hỏi--đáp nhiều bước
\cite{yang2018hotpotqa}, nhưng được điều chỉnh cho cấu trúc rule và event
của corpus pháp luật.

Benchmark không nhằm thay thế các bộ đánh giá pháp luật tổng quát như
LegalBench \cite{guha2023legalbench}. Nó được sử dụng như một testbed có
kiểm soát để kiểm tra pipeline trên chính loại cấu trúc nhân quả pháp luật
mà hệ thống được thiết kế để xử lý.

\subsection{Quy trình tạo câu hỏi}

Các câu hỏi được tạo từ causal path, rule và article có trong graph. Với một
đường hai bước:

\[
A \xrightarrow{r_1} B \xrightarrow{r_2} C,
\]

hệ thống có thể tạo các dạng câu hỏi sau:

\begin{description}
    \item[Forward multihop] Cho \(A\), yêu cầu xác định hệ quả cuối \(C\).

    \item[Reverse reasoning] Cho \(C\) hoặc chuỗi kết quả, yêu cầu xác định
    nguyên nhân ban đầu \(A\).

    \item[Bridge event] Cho \(A\) và \(C\), yêu cầu xác định mediator \(B\).

    \item[Causal chain explanation] Yêu cầu trình bày toàn bộ
    \(A \rightarrow B \rightarrow C\) và các căn cứ tương ứng.

    \item[Yes/no positive] Hỏi liệu quan hệ hoặc chuỗi được nêu có được graph
    hỗ trợ hay không.

    \item[Yes/no counterexample] Khẳng định một quan hệ trực tiếp
    \(A \rightarrow C\) trong khi graph chỉ hỗ trợ đường gián tiếp
    \(A \rightarrow B \rightarrow C\).

    \item[Branch reasoning] Yêu cầu tổng hợp nhiều outcome có chung prefix
    hoặc mediator.

    \item[Convergence reasoning] Yêu cầu tổng hợp nhiều điều kiện ban đầu
    cùng dẫn tới một suffix hoặc outcome.

    \item[Bridge-convergence fallback] Sử dụng cấu trúc hội tụ trong trường
    hợp không đủ điều kiện tạo một mẫu convergence hoàn chỉnh, nhưng vẫn có
    thể tạo câu hỏi về bridge hoặc path chính.
\end{description}

Các câu hỏi được tạo theo template có kiểm soát và sử dụng event name đọc
được thay vì event ID. Việc dùng template giúp xác định gold answer và gold
path một cách nhất quán, nhưng cũng có thể tạo ra các mẫu ngôn ngữ lặp lại.
Đây là một nguy cơ đối với khả năng tổng quát hóa và được xem là hạn chế của
benchmark.

\subsection{Phân bố loại câu hỏi}

Benchmark chứa tổng cộng \(250\) câu hỏi. Phân bố thực tế được trình bày
trong Bảng~\ref{tab:benchmark}.

\begin{table}[t]
    \centering
    \caption{Phân bố các loại câu hỏi trong benchmark.}
    \label{tab:benchmark}
    \begin{tabular}{lr}
        \hline
        \textbf{Loại câu hỏi} & \textbf{Số lượng} \\
        \hline
        Forward multihop & \(70\) \\
        Reverse reasoning & \(45\) \\
        Bridge event & \(40\) \\
        Causal chain explanation & \(35\) \\
        Yes/no positive & \(20\) \\
        Yes/no counterexample & \(20\) \\
        Branch reasoning & \(4\) \\
        Convergence reasoning & \(10\) \\
        Bridge-convergence fallback & \(6\) \\
        \hline
        \textbf{Tổng cộng} & \(\mathbf{250}\) \\
        \hline
    \end{tabular}
\end{table}

Bốn nhóm forward multihop, reverse reasoning, bridge event và causal chain
explanation chiếm phần lớn benchmark. Điều này phù hợp với mục tiêu chính là
đánh giá suy luận trên chuỗi hai bước. Các câu branch và convergence có số
lượng nhỏ hơn do graph chỉ chứa một số cấu trúc thỏa mãn yêu cầu nhiều path
cùng prefix hoặc suffix.

\subsection{Ground truth và schema đánh giá}

Mỗi benchmark sample có một mã định danh duy nhất và chứa hai nhóm trường:
trường mô tả câu hỏi và trường phục vụ evaluation.

Các trường mô tả chính gồm:

\begin{itemize}
    \item \texttt{id};
    \item \texttt{question};
    \item \texttt{question\_type};
    \item metadata về nguồn path hoặc cấu trúc tạo câu hỏi.
\end{itemize}

Object \texttt{evaluation} chứa ground truth cho các lớp đánh giá khác nhau:

\begin{description}
    \item[Retrieval ground truth]
    Bao gồm \texttt{gold\_rule\_ids},
    \texttt{gold\_article\_ids} và kích thước retrieval kỳ vọng.

    \item[Event ground truth]
    Bao gồm \texttt{gold\_event\_ids}, tên các event và final event.

    \item[Path ground truth]
    Bao gồm các bước của gold path theo thứ tự
    source event--target event.

    \item[Verification ground truth]
    Bao gồm \texttt{requires\_counterfactual},
    \texttt{gold\_decision} và target của quá trình xác minh.

    \item[Answer ground truth]
    Chứa câu trả lời tham chiếu phù hợp với question type.

    \item[Citation ground truth]
    Chứa các article ID được kỳ vọng hỗ trợ câu trả lời.

    \item[Difficulty metadata]
    Chứa độ khó của retrieval, reasoning, generation và một overall score.
\end{description}

Toàn bộ \(250\) câu hỏi có evaluation object, gold rule, gold event, gold
answer và gold citation. Có \(236\) câu hỏi có một gold path tuyến tính rõ
ràng. Các câu branch và convergence còn lại yêu cầu đánh giá trên nhiều path
hoặc một cấu trúc con thay vì một chuỗi tuyến tính duy nhất.

\subsection{Gold verification decision}

Phân bố gold decision gồm:

\begin{table}[t]
    \centering
    \caption{Phân bố gold verification decision.}
    \label{tab:gold-decision}
    \begin{tabular}{lr}
        \hline
        \textbf{Gold decision} & \textbf{Số lượng} \\
        \hline
        \texttt{SUPPORTED} & \(230\) \\
        \texttt{REJECT\_DIRECT\_CLAIM} & \(20\) \\
        \texttt{UNCERTAIN} & \(0\) \\
        \hline
        \textbf{Tổng cộng} & \(\mathbf{250}\) \\
        \hline
    \end{tabular}
\end{table}

Benchmark mất cân bằng mạnh về nhãn quyết định: \(92\%\) mẫu thuộc lớp
\texttt{SUPPORTED}, trong khi chỉ \(8\%\) thuộc lớp
\texttt{REJECT\_DIRECT\_CLAIM}. Vì vậy, verification accuracy không đủ để
đánh giá độc lập. Nghiên cứu báo cáo thêm Macro-F1, Balanced Accuracy và
confusion matrix.

Benchmark không chứa gold sample thuộc lớp \texttt{UNCERTAIN}. Do đó, khả
năng phát hiện thiếu bằng chứng hoặc trạng thái chưa xác định chưa được đánh
giá trực tiếp bằng một lớp gold riêng biệt.

\subsection{Counterfactual subset}

Benchmark đánh dấu \(20\) câu hỏi là yêu cầu xác minh counterfactual. Tuy
nhiên, việc kiểm tra nội dung cho thấy các mẫu này chủ yếu là
yes/no counterexample về quan hệ trực tiếp. Với đường:

\[
A \rightarrow B \rightarrow C,
\]

câu hỏi có thể khẳng định rằng \(A\) trực tiếp dẫn đến \(C\). Gold decision
khi đó là \texttt{REJECT\_DIRECT\_CLAIM} vì graph không chứa cạnh trực tiếp
\(A \rightarrow C\).

Các mẫu này đánh giá tốt khả năng phân biệt direct edge và indirect path,
nhưng không cung cấp ground truth đầy đủ cho phép can thiệp:

\[
do(B=\mathrm{FALSE}).
\]

Cụ thể, benchmark chưa gán nhãn chuyên gia cho:

\begin{itemize}
    \item factual state của mediator;
    \item factual state của outcome;
    \item counterfactual outcome sau intervention;
    \item các rule bị vô hiệu hóa;
    \item đường thay thế sau intervention;
    \item mediator có thực sự cần thiết trong mô hình pháp lý hay không.
\end{itemize}

Vì vậy, metric ``counterfactual verification accuracy'' trong nghiên cứu
hiện tại chủ yếu phản ánh khả năng xử lý direct-edge negative claims. Nó
không được diễn giải như bằng chứng rằng LegalSCM đã được xác nhận trên một
benchmark hard-intervention hoàn chỉnh.

\subsection{Độ khó của benchmark}

Benchmark gán difficulty metadata dựa trên đặc điểm retrieval, reasoning và
generation của từng mẫu. Trong tập hiện tại:

\begin{itemize}
    \item \(181\) câu được gán mức \texttt{medium};
    \item \(69\) câu được gán mức \texttt{hard};
    \item không có câu thuộc mức \texttt{easy}.
\end{itemize}

Các mẫu hard thường yêu cầu một trong các đặc điểm:

\begin{itemize}
    \item nhiều rule hoặc article hỗ trợ;
    \item nhiều candidate path có chung mediator hoặc outcome;
    \item cấu trúc branch hoặc convergence;
    \item phân biệt direct edge với indirect path;
    \item câu trả lời cần tổng hợp nhiều thành phần.
\end{itemize}

Difficulty label được sử dụng để phân tích kết quả theo nhóm, không được
truyền vào retriever, verifier hoặc answer generator.

\subsection{Kiểm soát rò rỉ nhãn đánh giá}

Do benchmark chứa \texttt{question\_type}, gold path và gold answer, cần
ngăn các trường này ảnh hưởng đến prediction. Pipeline áp dụng các nguyên
tắc sau:

\begin{enumerate}
    \item Step 3 chỉ nhận chuỗi \texttt{question}, không nhận evaluation
    object.

    \item Step 4 chỉ nhận retrieval result và query text được lưu trong
    artifact.

    \item Claim type được suy ra từ query text, không đọc
    \texttt{question\_type}.

    \item Step 5 suy ra answer intent từ query text và
    \texttt{query\_analysis} của Step 4.

    \item Gold rule, gold event, gold path, gold decision và gold answer chỉ
    được Step 7 sử dụng sau khi toàn bộ prediction đã được tạo.

    \item Batch runner có thể sao chép \texttt{question\_type} vào prediction
    để phục vụ phân tích theo nhóm, nhưng trường này không được chuyển vào
    các hàm retrieval, verification hoặc answer generation.
\end{enumerate}

Cơ chế này ngăn hệ thống sử dụng trực tiếp loại câu hỏi hoặc nhãn chuẩn để
chọn answer template. Đây là điều kiện cần để kết quả answer intent phản ánh
khả năng phân tích truy vấn thay vì rò rỉ từ benchmark.

\subsection{Tính toàn vẹn của paired experiment}

Hai run chính sử dụng cùng:

\begin{itemize}
    \item benchmark \(250\) câu;
    \item graph, causal memory, embedding và FAISS index;
    \item retriever và toàn bộ retrieval hyperparameter;
    \item query analyzer;
    \item Step 5 và extractive answer provider;
    \item evaluator và tập metric.
\end{itemize}

Khác biệt duy nhất là chế độ xác minh:

\begin{itemize}
    \item \texttt{structural\_scm};
    \item \texttt{path\_ablation}.
\end{itemize}

Cả hai run đều tạo đủ \(250/250\) prediction và không ghi nhận lỗi pipeline.
Việc giữ nguyên các thành phần còn lại cho phép xem đây là một paired
ablation hợp lệ ở mức implementation.

Các artifact cũ được tạo bằng runner hoặc Step 5 trước đó không được đưa vào
paired comparison. Cụ thể, chỉ các file có hậu tố
\texttt{\_v41\_step5\_v51} được sử dụng cho kết quả chính.

\subsection{Hạn chế và nguy cơ đối với tính hợp lệ của benchmark}

Benchmark hiện tại có một số hạn chế quan trọng.

\paragraph{Cùng nguồn graph.}

Câu hỏi và ground truth được tạo từ cùng graph mà pipeline sử dụng. Điều này
giúp kiểm soát cấu trúc suy luận nhưng có thể đánh giá cao khả năng khôi
phục các quan hệ đã được mã hóa. Benchmark chưa kiểm tra đầy đủ khả năng tổng
quát hóa sang một corpus hoặc graph độc lập.

\paragraph{Ảnh hưởng của template.}

Nhiều câu hỏi được tạo theo template. Các cụm từ như ``hệ quả cuối cùng'',
``sự kiện trung gian'', ``trực tiếp'' hoặc ``điều kiện ban đầu'' có thể giúp
query analyzer nhận diện answer intent. Hiệu năng trên câu hỏi tự nhiên có
cách diễn đạt đa dạng hơn có thể thấp hơn.

\paragraph{Thiếu thẩm định chuyên gia.}

Benchmark chưa được chuyên gia pháp luật rà soát toàn bộ ở mức từng câu,
từng path và từng answer. Gold data phản ánh cấu trúc graph và rule đã chuẩn
hóa, không phải kết luận pháp lý chính thức.

\paragraph{Counterfactual coverage hạn chế.}

Counterfactual subset chỉ có \(20\) mẫu và chủ yếu kiểm tra direct-edge
negative claims. Benchmark chưa đủ nhạy để phân biệt hard intervention với
node-deletion reachability.

\paragraph{Gold path tuyến tính.}

Metric exact path dựa trên một chuỗi tuyến tính, trong khi branch và
convergence có thể yêu cầu một tập path hoặc subgraph. Do đó, exact path
bằng không ở các nhóm này không nhất thiết đồng nghĩa với câu trả lời sai.

\paragraph{Rule representation đơn giản hóa.}

Corpus chưa biểu diễn đầy đủ negation, exception, nhiều điều kiện đồng thời,
rule priority và xung đột pháp lý. Benchmark vì thế cũng chưa đánh giá các
năng lực này.

\paragraph{Phạm vi miền hẹp.}

Benchmark chỉ tập trung vào Bộ luật Hình sự và không đại diện cho các lĩnh
vực pháp luật khác như dân sự, hành chính, lao động hoặc tố tụng.

Do các hạn chế trên, kết quả thực nghiệm nên được hiểu là đánh giá nội bộ
trên một testbed có kiểm soát. Chúng không được sử dụng để kết luận rằng hệ
thống có thể đưa ra tư vấn pháp lý đáng tin cậy trong môi trường thực tế.
